
\subsubsection{H-E1: Existence of Cross-Dimensional Effects}

GPT-2 was fine-tuned on 100 TruthfulQA samples for 3 epochs using LoRA. Three replicates with different random seeds (42, 43, 44) generated variance in performance deltas. All three dimensions were evaluated pre- and post-intervention. Pairwise correlations were computed across the three replicates for each dimension pair. Statistical significance was assessed via two-tailed t-test on correlation coefficients.

\subsubsection{H-M1: Target Dimension Improvement}

Using the same three GPT-2 replicates, we tested whether fine-tuning improved truthfulness. Baseline TruthfulQA MC2 scores (evaluated on full 817 questions) were compared to post-intervention scores. Statistical significance was tested via paired t-test against null hypothesis μ(Δ) = 0. Directional consistency was measured as the percentage of replicates showing positive Δ.

\subsubsection{H-M2: Representation Propagation}

Internal activations were extracted from GPT-2 before and after LoRA intervention using TransformerLens. We analyzed 24 layers: 12 attention pattern layers (blocks.{0-11}.attn.hook_pattern) and 12 residual stream layers (blocks.{0-11}.hook_resid_post). For each layer, activations were extracted on 100 held-out samples and CKA similarity was computed between pre- and post-intervention representations. Change magnitude was measured as Δ = 1 - CKA. Correlation between layer-wise representation change and H-M1 performance delta was tested using Pearson correlation.

\subsubsection{H-M3: Selective Coupling}

H-E1 was scaled up with 500 TruthfulQA training samples (instead of 100) and full 3-epoch fine-tuning to generate stronger intervention effects. All three dimensions were evaluated pre- and post-intervention across three seeds. Pairwise correlations were computed for all dimension pairs (truthfulness-fairness, truthfulness-robustness, fairness-robustness). Statistical significance was assessed via t-test. Permutation tests (1,000 permutations) established null distributions by randomly shuffling dimension assignments.

\subsubsection{H-M4: Architectural Replication}

The H-M3 protocol was replicated across GPT-2, OPT-350M, and Pythia-410M. Each model underwent minimal LoRA perturbation (10 training steps instead of 3 epochs) to reduce computational cost. Five seeds per model (total: 15 runs) enabled robust correlation estimation per architecture. For each dimension pair, correlation direction was classified as positive (r > 0.3), negative (r < -0.3), or neutral (-0.3 ≤ r ≤ 0.3). Replication rate was computed as the percentage of models showing the majority direction.

\subsubsection{Training Configuration}

All experiments used consistent hyperparameters: learning rate 1e-4, AdamW optimizer, cosine learning rate schedule with 10\% warmup, batch size 4, gradient accumulation steps 2. LoRA configurations used rank 8, alpha 16, dropout 0.1, targeting c_attn modules. All experiments used single GPU training (CUDA device 0).

\subsubsection{Statistical Analysis}

Pearson correlation coefficients were computed for dimension pairs across replicates. Statistical significance was assessed using two-tailed t-tests. Permutation tests established null distributions for comparison. Statistical power was limited by small sample sizes (n = 3-5 per condition), providing approximately 50\% power to detect large effects (r ≥ 0.9) at α = 0.05.
